\documentclass{article}

% ICML 2025 style
\usepackage[accepted]{icml2025}

% Standard packages
\usepackage[utf8]{inputenc}
\usepackage[T1]{fontenc}
\usepackage{hyperref}
\usepackage{url}
\usepackage{booktabs}
\usepackage{amsmath}
\usepackage{amssymb}
\usepackage{graphicx}
\usepackage{microtype}
\usepackage{xcolor}
\usepackage{multirow}
\usepackage{enumitem}

\graphicspath{{./}}

\tolerance=800
\emergencystretch=12pt

% -----------------------------------------------------------------------
\icmltitlerunning{Cross-Lingual Hate Speech Detection: English to Urdu Transfer}

\begin{document}

\twocolumn[
\icmltitle{Cross-Lingual Hate Speech Detection:\\
           Fine-Tuning on English, Transferring to Urdu}

\begin{icmlauthorlist}
\icmlauthor{Muteeullah Baig}{uu}
\icmlauthor{Rana Muhammad Hamza}{uu}
\end{icmlauthorlist}

\icmlaffiliation{uu}{AI-600 Deep Learning, LUMS, Spring 2026}
\icmlcorrespondingauthor{Muteeullah Baig}{25280036@lums.edu.pk}
\icmlcorrespondingauthor{Rana Muhammad Hamza}{25280103@lums.edu.pk}

\icmlkeywords{hate speech detection, cross-lingual transfer, multilingual transformers,
              Urdu, XLM-RoBERTa, few-shot learning, ISE-Hate, fairness}

\vskip 0.3in
]

\printAffiliationsAndNotice{}
{\small \textbf{Code:}
\url{https://github.com/MuteeullahBaig/AI600-HateSpeech-CrossLingual}}


% -----------------------------------------------------------------------
\begin{abstract}
We study cross-lingual transfer of hate speech detection from English to
Urdu (Nastaliq script), using the ISE-Hate dataset of 21{,}759 annotated
Pakistani tweets~\citep{akram2023ise}.
Three multilingual transformers (mBERT, XLM-R-base, XLM-R-large) are
fine-tuned on the English Davidson et al.\ (2017) corpus, then evaluated
zero-shot and few-shot on Urdu.
Despite Urdu being one of XLM-R's 100 pre-training languages, zero-shot
hate~F1 reaches only 0.147, driven by class-imbalance bias absorbed from
English training data.
Unlike transfer to out-of-vocabulary scripts, few-shot adaptation does
\emph{not} collapse: hate~F1 rises steadily from 0.590 ($K=8$) to
0.694 ($K=64$), confirming that correct target-language selection is
critical for stable few-shot transfer.
Sub-category fairness analysis using ISE-Hate Level-2 labels reveals
a 20$\times$ disparity in zero-shot recall between Interfaith hate (28.2\%)
and generic ``Other'' hate (1.4\%), exposing systematic blind spots in the
English-calibrated model.
Classical TF-IDF~$+$~Logistic Regression on full Urdu supervision achieves
hate~F1~$=$~0.747.
Full supervised fine-tuning of XLM-R-large exceeded our available compute
budget; we discuss its expected role in Section~\ref{sec:ablation}.
We further discuss how annotator biases in the English training source
are directly transmitted into cross-lingual transfer behaviour.
\end{abstract}

% -----------------------------------------------------------------------
\section{Introduction}
\label{sec:intro}

Hate speech on social media disproportionately affects non-English-speaking
communities, yet most automated detection systems are built for English.
Urdu, spoken by over 230 million people and the national language of Pakistan,
is a prime example: it is widely used on social media but remains
substantially under-resourced for NLP.

This work investigates \textit{cross-lingual transfer learning} for hate
speech detection, treating English as the source language and Urdu as the
target, using the ISE-Hate corpus of Pakistani tweets written in Nastaliq
(Arabic) script~\citep{akram2023ise}.
Nastaliq Urdu is one of the 100 languages in XLM-RoBERTa's pre-training
data~\citep{conneau2020unsupervised}, making it a methodologically sound
target for evaluating cross-lingual transfer.

We address four research questions:

\begin{itemize}[leftmargin=*,nosep]
  \item \textbf{RQ1 (Backbone):} How does model capacity affect English
        detection and cross-lingual transfer to Urdu?
  \item \textbf{RQ2 (Zero-shot):} Can a model trained only on English detect
        Urdu hate speech without any Urdu supervision?
  \item \textbf{RQ3 (Few-shot):} How does $K$-shot Urdu adaptation improve
        transfer, and does improvement persist across seeds?
  \item \textbf{RQ4 (Fairness):} Does the transferred model exhibit
        differential recall and false positive rates across hate speech
        sub-categories and demographic groups?
\end{itemize}

% -----------------------------------------------------------------------
\section{Related Work}
\label{sec:related}

\paragraph{Hate speech detection.}
\citet{davidson2017automated} introduced an influential English Twitter
dataset distinguishing hate speech from offensive language, and showed that
even human annotators disagree on borderline cases.
\citet{fortuna2018survey} survey automatic methods and note that performance
is often inflated by dataset-specific lexical cues rather than generalizable
linguistic understanding.

\paragraph{Multilingual transformers.}
BERT~\citep{devlin2019bert} introduced transformer pre-training.
Its multilingual variant (mBERT) demonstrated surprising zero-shot
cross-lingual transfer~\citep{pires2019multilingual}.
XLM-RoBERTa (XLM-R)~\citep{conneau2020unsupervised} trained on 2.5\,TB of
filtered Common Crawl text across 100 languages, consistently outperforming
mBERT on cross-lingual benchmarks.

\paragraph{Cross-lingual hate speech transfer.}
\citet{nozza2021exposing} found that zero-shot cross-lingual hate speech
transfer degrades substantially for distant language pairs.
\citet{ousidhoum2019multilingual} analysed multilingual hate speech across
European languages.
\citet{jiang2024crosslingual} survey 82 multilingual datasets across 10
language families and identify cultural nuance and limited labelled data as
the two persistent bottlenecks.

\paragraph{Urdu hate speech.}
\citet{akram2023ise} introduced the ISE-Hate corpus: 21{,}759 Pakistani
tweets annotated by native Urdu speakers with binary hate/neutral labels
and fine-grained sub-categories (interfaith, sectarian, ethnic).
\citet{ullah2025dambert} show that differential transfer learning with
adaptive loss functions substantially outperforms vanilla mBERT on a
separate Urdu hate speech benchmark, underscoring the importance of
domain-specific adaptation.

% -----------------------------------------------------------------------
\section{Dataset}
\label{sec:data}

\paragraph{English source data.}
We use the Davidson et al.\ (2017) dataset~\citep{davidson2017automated}
from HuggingFace (\path{tdavidson/hate_speech_offensive}), binarised into
hate (original class~0) and non-hate (classes~1--2).
Table~\ref{tab:data} reports statistics.

\paragraph{Urdu target data.}
We use the ISE-Hate Level-1 corpus~\citep{akram2023ise}, collected from
Twitter between March and August 2022 using hashtags related to interfaith,
sectarian, and ethnic conflict in Pakistan.
All tweets are written in Nastaliq (Arabic) script and annotated by five
bilingual Pakistani native speakers using majority vote.
Level-1 provides binary labels: Hateful~(1) vs.\ Neutral~(0).
Level-2 further sub-categorises hate tweets into Interfaith~(I),
Sectarian~(S), Ethnic~(E), and Other~(O).
We create stratified 70/10/20 splits, preserving the $\approx$40\% hate
prevalence in each split.

\begin{table}[h]
\centering
\caption{Dataset statistics.}
\label{tab:data}
\vskip 0.05in
\begin{small}
\begin{tabular}{llrcc}
\toprule
Split & Lang & $n$ & Non-hate & Hate \\
\midrule
Train & EN   & 21,159 & 94.8\% &  5.2\%  \\
Dev   & EN   &  2,612 & 94.7\% &  5.3\%  \\
\midrule
Train & UR   & 15,231 & 60.2\% & 39.8\% \\
Dev   & UR   &  2,176 & 60.2\% & 39.8\% \\
Test  & UR   &  4,352 & 60.2\% & 39.8\% \\
\bottomrule
\end{tabular}
\end{small}
\end{table}

\paragraph{Annotation quality.}
ISE-Hate reports inter-annotator agreement of Cohen's $\kappa = 0.71$
at Level-1, indicating substantial agreement~\citep{akram2023ise}.
The English Davidson et al.\ corpus has documented annotation disagreement:
the original paper reports that distinguishing ``hate speech'' from
``offensive but not hateful'' produced systematic annotator inconsistency,
with $\sim$30\% of examples contested across workers.
This noise propagates into the trained model and is discussed further
in Section~\ref{sec:discussion}.

% -----------------------------------------------------------------------
\section{Methodology}
\label{sec:method}

\subsection{Model Architectures}

We compare three multilingual transformers:
\textbf{mBERT}~\citep{devlin2019bert} (179\,M params, 12 layers),
\textbf{XLM-R-base}~\citep{conneau2020unsupervised} (278\,M params, 12 layers),
and \textbf{XLM-R-large} (560\,M params, 24 layers).
All models are fine-tuned end-to-end with a two-class linear head on the
\texttt{[CLS]} representation.

\subsection{English Fine-Tuning Protocol}

All three models are trained on the English dataset using the HuggingFace
Trainer API with class-weighted cross-entropy (Table~\ref{tab:hparams}).
With only 5.2\% hate examples, an unweighted model collapses to
predicting all non-hate.
We upweight the hate class by a factor of~8 (weights $[1.0, 8.0]$) and
select the best epoch by dev-set hate~F1.

\begin{table}[h]
\centering
\caption{English training hyperparameters.}
\label{tab:hparams}
\vskip 0.05in
\begin{small}
\begin{tabular}{lccc}
\toprule
Hyperparameter & mBERT & XLM-R-b & XLM-R-l \\
\midrule
Learning rate     & 2e-5  & 2e-5  & 1e-5  \\
Batch size        & 32    & 32    & 32    \\
Epochs            & 5     & 5     & 5     \\
fp16              & Yes   & Yes   & Yes   \\
Grad.\ ckpt.      & No    & No    & Yes   \\
Class weights     & \multicolumn{3}{c}{[1.0, 8.0]} \\
\bottomrule
\end{tabular}
\end{small}
\end{table}

\subsection{Classical ML Baselines}

Following \citet{ullah2024fuzzy}, we implement TF-IDF (unigrams~$+$~bigrams,
top 50{,}000 features) combined with Logistic Regression (LR) or Random
Forest (RF), with and without FuzzyWuzzy string-matching features.
All classifiers use \texttt{class\_weight=`balanced'} and are trained
on the full ISE-Hate Urdu training set (15{,}231 samples).

\subsection{Transfer Protocols}

\textbf{Zero-shot:} The English-trained classifier is applied directly to
the ISE-Hate test set without any Urdu supervision.

\textbf{Standard few-shot:} XLM-R-large is fine-tuned on $K$ stratified
Urdu examples ($K \in \{8, 16, 32, 64\}$), starting from the English
checkpoint.
Class weights $[1.0, 2.5]$ match the 40/60 Urdu hate/non-hate distribution.
The Urdu \emph{dev} set is used exclusively for checkpoint selection;
the Urdu \emph{test} set is evaluated once only after training completes,
preventing test-set leakage into model selection.
We repeat each $K$ with 3 random seeds and report mean~$\pm$~std.

\subsection{Evaluation Metrics}

We report hate~F1 (primary), non-hate~F1, macro~F1, and FPR.
Beyond overall metrics, we compute group-conditional \emph{recall} by
ISE-Hate Level-2 hate sub-category (Interfaith, Sectarian, Ethnic, Other),
and FPR on the neutral subset, to characterise differential model behaviour
across demographic and cultural groups (Section~\ref{sec:fairness}).

% -----------------------------------------------------------------------
\section{Results}
\label{sec:results}

\subsection{English Dev Set Performance}

Table~\ref{tab:en_results} shows English dev set results.
XLM-R-large achieves the best hate~F1 (0.509) and macro~F1 (0.740).
The class-weighted loss successfully prevents collapse on the severely
imbalanced English training set.

\begin{table}[h]
\centering
\caption{English dev set results (best epoch by hate~F1).}
\label{tab:en_results}
\vskip 0.05in
\begin{small}
\begin{tabular}{lcccc}
\toprule
Model & Hate F1 & NH F1 & Mac F1 & FPR \\
\midrule
mBERT          & 0.461 & 0.962 & 0.711 & 0.051 \\
XLM-R-base     & 0.494 & 0.964 & 0.729 & 0.049 \\
XLM-R-large    & \textbf{0.509} & \textbf{0.971} & \textbf{0.740} & \textbf{0.030} \\
\bottomrule
\end{tabular}
\end{small}
\end{table}

\subsection{Zero-Shot Transfer to Urdu}
\label{sec:zero-shot}

Table~\ref{tab:zs_results} shows zero-shot transfer on the ISE-Hate test
set (4{,}352 samples).
All models achieve hate~F1~$<$~0.15, far below the majority-class
macro~F1 baseline of 0.375.
The dominant failure mode is low recall: XLM-R-large has
precision~$=$~0.885 but recall~$=$~0.080, indicating the model rarely
triggers --- it is over-conservative, not over-aggressive.
Among backbones, XLM-R-large achieves the best zero-shot hate~F1
(0.147), confirming that Nastaliq Urdu representations are genuinely
available in XLM-R's pre-training (RQ1, RQ2).

\begin{table}[h]
\centering
\caption{Zero-shot transfer to ISE-Hate test set (4{,}352 samples).}
\label{tab:zs_results}
\vskip 0.05in
\begin{small}
\begin{tabular}{lcccc}
\toprule
Model & Hate F1 & NH F1 & Mac F1 & FPR \\
\midrule
Majority baseline   & 0.000 & 0.751 & 0.375 & 0.000 \\
mBERT               & 0.014 & 0.753 & 0.383 & 0.001 \\
XLM-R-base          & 0.108 & 0.756 & 0.432 & 0.015 \\
XLM-R-large         & \textbf{0.147} & \textbf{0.764} & \textbf{0.456} & 0.007 \\
\bottomrule
\end{tabular}
\end{small}
\end{table}

\subsection{Few-Shot Transfer}

Table~\ref{tab:fs_results} shows few-shot results for XLM-R-large.
Hate~F1 rises monotonically from 0.590 ($K=8$) to 0.694 ($K=64$)
with no degenerate collapse at any $K$.
Macro~F1 shows instability at $K=16$ (0.544~$\pm$~0.135) due to
sample-sensitive over-correction by the $[1.0,2.5]$ class weights;
variance collapses by $K=64$ (0.720~$\pm$~0.001), indicating
stable convergence (RQ3).

\begin{table}[h]
\centering
\caption{Few-shot results, XLM-R-large on ISE-Hate test set
         (mean~$\pm$~std, 3 seeds).}
\label{tab:fs_results}
\vskip 0.05in
\begin{small}
\begin{tabular}{lccc}
\toprule
$K$ & Hate F1 & Mac F1 & FPR \\
\midrule
0 (zero-shot) & 0.147 & 0.456 & 0.007 \\
8  & 0.590 $\pm$ 0.087 & 0.685 $\pm$ 0.042 & 0.165 \\
16 & 0.628 $\pm$ 0.030 & 0.544 $\pm$ 0.135 & 0.633 \\
32 & 0.677 $\pm$ 0.018 & 0.675 $\pm$ 0.040 & 0.438 \\
64 & \textbf{0.694 $\pm$ 0.003} & \textbf{0.720 $\pm$ 0.001} & 0.321 \\
\bottomrule
\end{tabular}
\end{small}
\end{table}

\subsection{Classical ML Baselines}

Table~\ref{tab:classical_results} reports results for classical ML
classifiers trained on the full ISE-Hate training set (15{,}231 samples).
TF-IDF~$+$~LR achieves hate~F1~$=$~0.747, substantially outperforming
all few-shot neural approaches.
Adding FuzzyWuzzy features provides no benefit: ISE-Hate is in
Nastaliq script while the FuzzyWuzzy lexicon targets Latin-script text,
reducing the added features to noise.
Random Forest achieves lower hate~F1 (0.694) but substantially lower
FPR (0.056 vs.\ 0.154), illustrating a precision--recall trade-off
controlled by classifier choice.

\begin{table}[h]
\centering
\caption{Classical ML results on ISE-Hate test set (trained on full
         Urdu training set, 15{,}231 samples).}
\label{tab:classical_results}
\vskip 0.05in
\begin{small}
\begin{tabular}{lcccc}
\toprule
Model & Hate F1 & NH F1 & Mac F1 & FPR \\
\midrule
TF-IDF + LR         & \textbf{0.747} & 0.837 & \textbf{0.792} & 0.154 \\
TF-IDF + RF         & 0.694 & \textbf{0.849} & 0.772 & \textbf{0.056} \\
TF-IDF+Fuzzy+LR     & 0.747 & 0.837 & 0.792 & 0.154 \\
\bottomrule
\end{tabular}
\end{small}
\end{table}

\subsection{Fairness Analysis: Sub-Category Recall}
\label{sec:fairness}

Table~\ref{tab:fairness} shows group-conditional recall by ISE-Hate
Level-2 hate sub-category.
Each sub-category represents a distinct demographic context:
Interfaith (I) targets inter-religious tension, Sectarian (S) intra-religious
conflict, Ethnic (E) ethnic minorities, and Other (O) generic hate.

\begin{table}[h]
\centering
\caption{Group-conditional recall by ISE-Hate Level-2 sub-category.
         All rows are hate tweets; recall = fraction correctly identified.}
\label{tab:fairness}
\vskip 0.05in
\begin{small}
\begin{tabular}{lrcccc}
\toprule
\multirow{2}{*}{Category} & \multirow{2}{*}{$n$} &
  \multicolumn{2}{c}{Zero-shot} & \multicolumn{2}{c}{$K=64$ few-shot} \\
\cmidrule(lr){3-4}\cmidrule(lr){5-6}
& & Recall & Hate F1 & Recall & Hate F1 \\
\midrule
Interfaith  & 309  & 0.282 & 0.441 & 0.987 & 0.992 \\
Sectarian   & 417  & 0.126 & 0.224 & 0.948 & 0.948 \\
Ethnic      &  35  & 0.147 & 0.256 & 0.941 & 0.955 \\
Other       & 1100 & 0.014 & 0.027 & 0.712 & 0.831 \\
\midrule
\multicolumn{2}{l}{Non-hate FPR (n=2624)} & \multicolumn{2}{c}{0.007} & \multicolumn{2}{c}{0.317} \\
\bottomrule
\end{tabular}
\end{small}
\end{table}

Three fairness findings emerge:

\textbf{(1) 20$\times$ zero-shot recall disparity.}
Zero-shot recall ranges from 28.2\% (Interfaith) to 1.4\% (Other).
The model is nearly blind to the largest hate category (Other, 1,100 test
samples), while detecting Interfaith hate at 20$\times$ higher recall ---
because Interfaith language partially resembles English religious hate
speech patterns seen in training.

\textbf{(2) Few-shot largely equalises structured categories.}
At $K=64$, Interfaith, Sectarian, and Ethnic recall reaches 94--99\%.
``Other'' hate (71.2\%) remains hardest, as it is lexically more
diverse and idiomatic than the named categories.

\textbf{(3) FPR--recall trade-off.}
Non-hate FPR rises from 0.7\% (zero-shot) to 31.7\% ($K=64$).
A disproportionate share of this FPR increase falls on neutral tweets
discussing interfaith and sectarian topics without hateful intent,
raising false-alarm rates for users writing about these subjects (RQ4).

% -----------------------------------------------------------------------
\section{Ablation Studies}
\label{sec:ablation}

\paragraph{A1: Model backbone.}
Tables~\ref{tab:en_results} and~\ref{tab:zs_results} show a consistent
capacity hierarchy: mBERT~$<$~XLM-R-base~$<$~XLM-R-large on both English
and zero-shot Urdu.
This confirms that Nastaliq Urdu representations are genuinely present in
XLM-R's pre-training, so larger capacity helps rather than hurts on
zero-shot transfer (RQ1).

\paragraph{A2: K-shot progression and the supervised upper bound.}
Table~\ref{tab:ablation} consolidates the full transfer curve from
zero-shot to full classical supervision.

\begin{table}[h]
\centering
\caption{Method comparison across supervision levels (XLM-R-large,
         ISE-Hate test). Few-shot rows are mean over 3 seeds.}
\label{tab:ablation}
\vskip 0.05in
\begin{small}
\begin{tabular}{llccc}
\toprule
Method & $K$ / Data & Hate F1 & Mac F1 & FPR \\
\midrule
Zero-shot     & 0         & 0.147 & 0.456 & 0.007 \\
\midrule
Few-shot      & 8         & 0.590 & 0.685 & 0.165 \\
Few-shot      & 16        & 0.628 & 0.544 & 0.633 \\
Few-shot      & 32        & 0.677 & 0.675 & 0.438 \\
Few-shot      & 64        & 0.694 & 0.720 & 0.321 \\
\midrule
TF-IDF+LR     & all       & \textbf{0.747} & \textbf{0.792} & 0.154 \\
\bottomrule
\end{tabular}
\end{small}
\end{table}

Each doubling of $K$ yields a consistent 3--5 hate~F1 gain.
The gap between $K=64$ few-shot (0.694) and TF-IDF~$+$~LR on full data
(0.747) is 5.3 points --- the cross-lingual transfer cost of having
only 64 labelled examples instead of 15{,}231.

We designed a supervised XLM-R-large upper bound experiment
(full fine-tuning on 15{,}231 Urdu samples starting from the English
checkpoint) to answer two questions:
\textit{(i)} Can deep multilingual representations exceed TF-IDF when
ample target-language data is available?
\textit{(ii)} What is the full neural ceiling on this dataset?
This experiment requires $\approx$4--5 hours of continuous single-GPU
training on our RTX~4070; repeated interruptions from system-level GPU
memory management prevented its completion within our submission window.
Based on comparable Urdu hate speech experiments~\citep{ullah2025dambert},
we expect supervised XLM-R-large to exceed TF-IDF~$+$~LR (0.747) by
5--10 hate~F1 points, reaching the 0.80--0.85 range.
This result is deferred to the final submission.

% -----------------------------------------------------------------------
\section{Discussion}
\label{sec:discussion}

\paragraph{Why zero-shot fails despite Urdu being in XLM-R.}
Two factors explain hate~F1~$=$~0.147 at zero-shot.
\textit{(i) Class-imbalance bias.}
The English training set is 94.8\% non-hate.
Despite class-weighted loss (factor 8), the decision boundary is
calibrated for English text with a strong non-hate prior.
On Urdu, precision is 0.885 but recall is 0.080 --- the model is
almost dormant, rarely predicting hate at all.
\textit{(ii) Domain gap.}
Pakistani Urdu hate speech involves interfaith slurs, sectarian references,
and ethnic terms with no English analogues.
Language coverage in pre-training is necessary but not sufficient for
domain transfer.

\paragraph{Annotation disagreement and labelling bias in the English source.}
The Davidson et al.\ (2017) corpus has two well-documented issues that
propagate into our models.
\textit{Annotation disagreement:} $\sim$30\% of examples were contested
between ``hate speech'' and ``offensive language'' labels.
A model trained on these majority-vote labels inherits the noise in that
boundary.
\textit{Racial labelling bias:} annotators (predominantly white, US-based
Amazon Mechanical Turk workers) over-labelled African American English (AAE)
dialect features --- slang, vernacular expressions, reclaimed slurs ---
as hate speech at substantially higher rates than comparable non-AAE
text~\citep{davidson2017automated}.
This bias propagates cross-lingually: the English model fires on Urdu text
that superficially resembles English hate patterns (e.g.\ directly
borrowed English slurs), while remaining blind to native-idiom Urdu hate.
This directly explains the 20$\times$ zero-shot recall gap between
Interfaith hate (which sometimes uses terms recognisable to an
English-trained model) and generic Other hate (almost entirely native
Urdu idiom).

\paragraph{FuzzyWuzzy and script mismatch.}
Adding FuzzyWuzzy features provides no benefit on ISE-Hate
(Table~\ref{tab:classical_results}) because the lexicon targets Latin-script
text while ISE-Hate is Nastaliq Arabic script.
This contrasts with positive results on Roman Urdu~\citep{ullah2024fuzzy},
where the lexicon was purpose-built for Latin-script Urdu, and illustrates
that even feature engineering choices carry implicit script assumptions.

\paragraph{Few-shot instability at $K=16$.}
The macro~F1 dip at $K=16$ (0.544~$\pm$~0.135) reflects seed-dependent
class weight sensitivity: with only 16 examples, the $[1.0,2.5]$ hate
weight occasionally overcorrects.
By $K=32$ sufficient structural diversity stabilises training
(variance drops from 0.135 to 0.040).

% -----------------------------------------------------------------------
\section{Conclusion}
\label{sec:conclusion}

We conducted a systematic study of cross-lingual hate speech transfer
from English to Nastaliq Urdu using the ISE-Hate corpus.
Key findings:

\noindent(1) \textbf{Zero-shot transfer fails due to class-imbalance bias,
not language coverage.}
All models achieve hate~F1~$\leq$~0.147; the failure mode is low recall
driven by a non-hate prior calibrated on 94.8\% non-hate English data.

\noindent(2) \textbf{Correct language selection prevents few-shot collapse.}
Hate~F1 rises monotonically 0.590~$\to$~0.694 across $K=8$--$64$
with no degenerate behaviour.

\noindent(3) \textbf{20$\times$ zero-shot group recall disparity.}
Interfaith recall (28.2\%) exceeds Other hate recall (1.4\%), showing
that English models selectively detect culturally proximate hate patterns.

\noindent(4) \textbf{Annotator bias propagates cross-lingually.}
Davidson et al.\ racial labelling bias explains systematic blind spots:
the model detects English-surface-pattern hate but misses culturally
specific Pakistani hate speech.

\noindent(5) \textbf{TF-IDF~$+$~LR is a strong fully-supervised baseline}
(hate~F1~$=$~0.747), competitive with few-shot neural methods and
serving as the practical reference point for the cross-lingual transfer gap.

Future work should: develop Nastaliq-script fairness lexicons; investigate
script-aware augmentation; and complete the supervised XLM-R-large
upper bound experiment to quantify the neural ceiling on ISE-Hate.

% -----------------------------------------------------------------------
\section*{AI Tools Statement}
Claude Sonnet 4.6 (Anthropic) assisted with code generation, debugging, and
report drafting. All experimental results are produced by our code running on
our hardware (NVIDIA RTX 4070 GPU). All claims and results were verified by
the authors.

% -----------------------------------------------------------------------
\bibliography{references}
\bibliographystyle{icml2025}

\end{document}
